\section{Capítulo~\ref{sec:nora}. \nt{NORA}}

La organización del sistema \nt{NORA}, sus dos modos, el vínculo con la pupila (no solo a través de \nt{NORA}), el aumento de la relación señal/fondo y la U invertida en la conducta están medidos directamente; la ganancia multiplicativa $g$ es un modelo; que la ganancia conmute el modo de elección y actúe como temperatura inversa son deducciones del capítulo; el beneficio de ajustar la ganancia es un cálculo con el modelo del libro; “\nt{NORA} es la perilla de la exploración” y “\nt{NORA} es una interrupción” son hipótesis.

{\footnotesize
\setlength{\tabcolsep}{3pt}\renewcommand{\arraystretch}{1.15}
\begin{longtable}{>{\raggedright}p{0.5cm}>{\raggedright}p{4.0cm}>{\raggedright}p{5.6cm}>{\raggedright}p{2.3cm}>{\raggedright\arraybackslash}p{1.7cm}}
\toprule
N.º & Proposición & Experimento y qué se midió & Fuente & Estado \\
\midrule\endfirsthead
\toprule
N.º & Proposición & Experimento y qué se midió & Fuente & Estado \\
\midrule\endhead
1 & En el ser humano hay unas decenas de miles de neuronas \nt{NORA}, casi todas en un solo grupo & Cortes seriados del tronco encefálico humano, tinción de tirosina hidroxilasa: $53\,900$ células en el locus coeruleus y otras $6\,260$ en el subnúcleo vecino. & \cite{Baker1989LC} & medido \\
2 & Las salidas se extienden por casi todo el cerebro, pero partes del grupo atienden en parte a regiones distintas & Marcaje viral-genético en ratones: las neuronas \nt{NORA} del locus coeruleus reciben entradas de amplias regiones del cerebro y envían salidas a amplias regiones del cerebro y de la médula espinal. En ratas: el grupo que va a la amígdala refuerza el aprendizaje del miedo, y un grupo aparte que va a la corteza prefrontal, su extinción. & \cite{Schwarz2015LC, Uematsu2017LC, Poe2020} & medido \\
3 & Modo de fondo: espigas regulares con una frecuencia de fracciones de espiga a unas pocas por segundo, que cambia con la vigilia & Ratas, registro de neuronas del locus coeruleus a lo largo del ciclo sueño-vigilia: la frecuencia es máxima en vigilia, menor en el sueño lento y casi nula en el sueño REM; los cambios de frecuencia se adelantan al cambio de estado. & \cite{AstonJonesBloom1981} & medido \\
4 & Modo fásico: una ráfaga breve ante un suceso importante para la tarea, aproximadamente en el momento de decidir & Monos, tarea de detectar un blanco poco frecuente: todas las neuronas del locus coeruleus responden con una ráfaga solo al blanco, $\sim90\ms$ después de él y $\sim200\ms$ antes de la respuesta manual; cuando el desempeño es malo, las ráfagas son más débiles. En una tarea de elección entre dos, la ráfaga está más ligada a la respuesta que al estímulo, y aparece tanto en respuestas correctas como erróneas. & \cite{AstonJones1994vigilance, Clayton2004LC} & medido \\
5 & El diámetro de la pupila es un punto de control impreciso de \nt{NORA}: sigue a \nt{NORA}, a \nt{ACET} y a otras regiones & Monos: las espigas del locus coeruleus, hasta espigas sueltas de una célula, predicen la dilatación de la pupila; hay un vínculo parecido pero más tardío en los colículos y la corteza cingulada. Ratones: las fluctuaciones de la pupila siguen la actividad de las salidas tanto noradrenérgicas como colinérgicas en la corteza. & \cite{Joshi2016, Reimer2016} & medido \\
6 & La noradrenalina suprime la actividad de fondo más que la evocada; la ganancia multiplicativa~\eqref{eq:gain} es un modelo que reproduce el efecto & Monos despiertos, corteza auditiva, iontoforesis de noradrenalina: la actividad de fondo de las neuronas se suprime más que la respuesta a las vocalizaciones de congéneres; la relación señal-ruido crece. La sigmoide multiplicativa~\eqref{eq:gain} procede de un modelo de red; la multiplicación literal de la pendiente no se ha medido. & \cite{Foote1975NE, ServanSchreiber1990} & medido; $g$, modelo \\
7 & La ganancia aumenta $d'$ solo frente al ruido posterior al amplificador~\eqref{eq:gainsnr} (proposición~\ref{prop:gainnoise}) & Deducción en el capítulo; en el modelo de red, la ganancia mejora la detección de la señal. No hay experimento que separe el ruido anterior y posterior al amplificador y compruebe esta diferencia. & deducción en el capítulo; \cite{ServanSchreiber1990} & teoría \\
8 & La frontera $g\beta=1$ conmuta la red de elección de “delibera” a “ha decidido”~\eqref{eq:wtagain} (proposición~\ref{prop:gainwta}, fig.~\ref{fig:pitchfork}) & Análisis lineal de dos grupos que se inhiben mutuamente; deducción en el capítulo. No hay medidas directas de este umbral en el cerebro. & deducción en el capítulo & teoría \\
9 & La ganancia es la temperatura inversa de la elección~\eqref{eq:softmax} & Con ruido logístico después del amplificador, la probabilidad de elección es un softmax con $T=\sigma/g$; deducción en el capítulo. & deducción en el capítulo & teoría \\
10 & \nt{NORA} fija cuánto explorar: una actividad de fondo alta es una $g$ efectiva pequeña (exploración); la ráfaga fásica en el momento de decidir, una grande (decisión) & Humanos: antes de una elección al azar la pupila basal es más ancha que antes de elegir la mejor opción conocida; las personas con pupila más ancha tienden más a explorar. Ratas: el paso al modo de elección aleatoria lo controla la entrada de \nt{NORA} a la corteza cingulada anterior. Que la ráfaga aumente la $g$ efectiva no se ha medido directamente. & \cite{JepmaNieuwenhuis2011, Tervo2014stochastic, AstonJones2005} & hipótesis \\
11 & La recompensa depende de $g$ como una U invertida; la mejor $g$ es menor en un mundo que cambia deprisa (proposición~\ref{prop:explore}, fig.~\ref{fig:invertedU}) & \texttt{models/nora\_explore.py}, $60$ ejecuciones de $4000$ intentos. Un cambio cada $1000$ intentos: mejor $g=16$, fracción $0.976$; cada $20$: $g=8$, fracción $0.886$; con $g=0.5$, $0.77$. El recálculo coincidió con el capítulo. & \texttt{models/\allowbreak nora\_\allowbreak explore.py} & modelo \\
12 & U invertida en la conducta: el mejor desempeño, con actividad de fondo moderada y ráfagas nítidas & Monos, la misma tarea que en la fila~4: en los periodos de buen desempeño la frecuencia de fondo es moderada y las ráfagas ante el blanco son nítidas; con frecuencia de fondo alta no hay ráfagas y hay más falsas alarmas. Los cambios de la actividad de fondo y evocada siguen las fluctuaciones del desempeño. & \cite{Usher1999LC, AstonJones2005} & medido \\
13 & \nt{NORA} comunica la incertidumbre inesperada, y \nt{ACET}, la esperada & Teoría de inferencia bayesiana con dos clases de incertidumbre; en los trabajos citados no hay experimento directo que separe estas señales por neurotransmisor. & \cite{YuDayan2005} & hipótesis \\
14 & Tras un cambio brusco las personas aprenden más deprisa, y la pupila se dilata & Humanos, tarea de predicción con cambios repentinos: los cambios breves de la pupila siguen la estimación de la probabilidad de cambio y, junto con la pupila basal, predicen cuánto modifican la estimación los datos nuevos (tasa de aprendizaje); un cambio de pupila ajeno a la luz y a la tarea modifica esa tasa. Que aquí la pupila indique precisamente \nt{NORA} es una inferencia a partir de los datos indirectos de la fila~5. & \cite{Nassar2012} & medido \\
15 & La ráfaga fásica funciona como una interrupción: la red borra su estado & No hay experimento directo en que una ráfaga de \nt{NORA} borre el estado de la red; solo están medidas las ráfagas ante sucesos importantes (fila~4). & \cite{BouretSara2005} & hipótesis \\
16 & Ajustar $g$ o la tasa de aprendizaje según la sorpresa apenas mejora los mejores ajustes constantes & \texttt{models/nora\_explore.py}, mundo con épocas de $1000$ intentos (un cambio cada $1000$ y cada $20$). Mejor constante, $g=8$: $0.925$; ajuste de $g$: hasta $0.927$; ajuste de la tasa de aprendizaje: hasta $0.926$; tasa constante $0.4$: $0.928$. El recálculo coincidió con el capítulo. & \texttt{models/\allowbreak nora\_\allowbreak explore.py} & modelo \\
\bottomrule
\end{longtable}}
